\section{\texorpdfstring{\textbf{Capítulo 1: Projeto de
qualificação}}{Capítulo 1: Projeto de qualificação}}\label{capuxedtulo-1-projeto-de-qualificauxe7uxe3o}

\subsection{\texorpdfstring{\textbf{Introdução}}{Introdução}}\label{introduuxe7uxe3o}

O projeto a seguir propõe uma sociologia histórica e multiescalar do
processo de desindustrialização brasileiro. Como problema se indaga:
qual a relação entre transformações institucionais e o processo de
desindustrialização brasileiro? Seu objetivo geral é compreender como
instituições políticas e econômicas influenciaram esse processo em
diferentes escalas geográficas. Para isso, utiliza de métodos
comparativos e de rastreamento de processos. Empiricamente, analisa
materiais qualitativos, dos planos de desenvolvimento nacionais aos
projetos de consórcio intermunicipais, e quantitativos, com ênfase em
séries temporais elaboradas a partir das Contas Nacionais, das Contas
Regionais, do Produto Interno Bruto (PIB) Municipal e de pesquisas sobre
produção e emprego.

A hipótese central é que novos arranjos institucionais, frutos dos
ciclos de crescimento e crise da economia brasileira durante a segunda
metade do século XX, aceleraram e aprofundaram o processo de
desindustrialização brasileiro. Como consequência, houve em menor
crescimento econômico, menor sofisticação tecnológica dos setores
econômicos\footnote{Aqui a tecnologia é entendida tanto do ponto de
  vista de seus suportes materiais -- máquinas, patentes etc --, como do
  ponto de vista do acúmulo de conhecimento e seus suportes imateriais
  -- qualificação profissional, modelos organizacionais, técnicas de
  gestão etc.}, menor produtividade do trabalhador, piores condições de
trabalho e inserção mais precária do país na economia mundial.

Além deste capítulo, o projeto contém um capítulo específico para
discussão metodológica e outro para discussão sobre o fenômeno da
desindustrialização em sua escala global\footnote{Dada a natureza
  multiescalar da investigação, a revisão de literatura sobre os casos
  (os períodos históricos, estados e sub-regiões) será incorporada aos
  capítulos que discutem estes.}. Já o restante deste capítulo contém
quatro seções: 1- objeto e discussão da hipótese; 2 - justificativa e
breve histórico da desindustrialização brasileira; 3 - teoria geral; 4 -
os objetivos gerais e específicos e 5 - o cronograma proposto e a
estrutura da tese.

\subsection{\texorpdfstring{\textbf{1. Objeto e discussão da
hipótese}}{1. Objeto e discussão da hipótese}}\label{objeto-e-discussuxe3o-da-hipuxf3tese}

O objeto da investigação são os efeitos das transformações
institucionais políticas e econômicas sobre o processo de
desindustrialização brasileiro e seus desdobramentos. No país, a
desindustrialização, que pode numa simplificação ser entendida como a
redução contínua da participação do setor industrial na economia e no
estoque de empregos
\href{https://www.zotero.org/google-docs/?kiAalY}{(Tregenna, 2009)},
ocorreu de forma consistente e mais intensa a partir dos anos 1980
\href{https://www.zotero.org/google-docs/?B3wCNB}{(Oreiro e Feijó,
2010)}. É possível a partir das Contas Nacionais
\href{https://www.zotero.org/google-docs/?kIvHAO}{(IBGE, {[}s.d.{]})},
porém, identificar episódios prévios do fenômeno, como no período de
1961 até 1966. Tal recorte temporal, de aproximadamente 1961 até os dias
atuais, é marcado por significativos ciclos de crescimento e crise
econômica \href{https://www.zotero.org/google-docs/?LxtZo0}{(Souza e
Previdelli, 2022; Vianna e Villela, 2015)}, reincidentes crises
políticas (idem) e significativa variação da atividade industrial
\href{https://www.zotero.org/google-docs/?r3fjx6}{(Morceiro, 2021)}.

Ao longo desses ciclos, transformações institucionais significativas
foram implementadas, em grande parte em resposta a esses ciclos e por
influência de elites políticas e econômicas. Pode-se identificar
mudanças de regime político, de política econômica e tributária, criação
e privatização de empresas estatais, fortalecimento e enfraquecimento de
autarquias, mudanças na legislação trabalhista e tributária, mudanças de
regime cambial e nas regras do comércio externo, para citar algumas
\href{https://www.zotero.org/google-docs/?GdUIm4}{(Abreu, 2014; Souza e
Previdelli, 2022; Vianna e Villela, 2015)}. Alternou-se entre e se
combinou abertura comercial e protecionismo, liberalismo e
intervencionismo, autoritarismo e democracia.

O fenômeno também apresenta variações na escala estadual e sub-regional.
Sinal disso é que arranjos institucionais ligados a estados e
territórios -- subsídios fiscais, diferentes níveis salariais médios e
diferentes características da atuação sindical
\href{https://www.zotero.org/google-docs/?gbGa0G}{(Ramalho, 2004,
2015)}, por exemplo -- têm contribuído para a desconcentração industrial
brasileira \href{https://www.zotero.org/google-docs/?qvqIuo}{(CNI,
2021)}, que ocorre paralelamente à redução do peso da indústria no
agregado nacional. Em função disso, a pesquisa optou por comparar os
processos de desindustrialização nos estados com maiores parques
industriais do país, São Paulo e Rio de Janeiro. Já no nível da
sub-região, entendido aqui como um conjunto de municípios articulados
por certa identidade regional e projeto econômico comum, são
contrastados os território produtivo do ABC paulista e do Sul
Fluminense\footnote{Ambos os territórios serão descritos mais adiante. O
  termo ``ABC paulista'' contempla os municípios de Santo André, São
  Bernardo do Campo, São Caetano do Sul, Diadema, Mauá, Ribeirão Pires e
  Rio Grande da Serra. É considerado parte da Região Metropolitana de
  São Paulo \href{https://www.zotero.org/google-docs/?Zpuw90}{(Governo
  do Estado de São Paulo, {[}s.d.{]})}, mas possui uma identidade e
  articulação política próprias, ilustrada no Consórcio Intermunicipal
  do ABC \href{https://www.zotero.org/google-docs/?JmJ6HS}{(Consórcio
  ABC, {[}s.d.{]})}. Por sua vez, o Sul fluminense consiste numa
  denominação que tem sido aplicada para um conjunto de municípios que
  hoje são parte da Região Geográfica Intermediária de Volta
  Redonda-Barra Mansa
  \href{https://www.zotero.org/google-docs/?A380nY}{(IBGE, 2017)}.
  Desses, os municípios de Itatiaia, Porto Real, Quatis e Resende
  constituem a Região Geográfica Imediata de Resende, onde se concentram
  as citadas atividades da indústria automotiva
  \href{https://www.zotero.org/google-docs/?9aZc1J}{(Ramalho e Santos,
  2022)}. Acrescidos dos 8 municípios da região Geográfica imediata de
  Volta Redonda-Barra Mansa, em especial de Volta Redonda, antigo pólo
  siderúrgico nacional, constituem o foco da análise quando se fala em
  ``Sul fluminense''.}, exemplos de estratégias de reconversão
industrial e re-industrialização e entre os quais estabelece-se uma
dinâmica de realocação produtiva ligada ao setor automotivo
\href{https://www.zotero.org/google-docs/?w7wz1h}{(Rodrigues e Ramalho,
2007)}.

Em função das discussões acima, se retorna à hipótese já apresentada de
que as transformações institucionais brasileiras impactaram o processo
de desindustrialização brasileiro de forma negativa, acelerando seu
ritmo e aprofundando seus impactos. Isto se daria pois tais arranjos
institucionais, ao responder às diferentes conjunturas políticas,
econômicas e sociais, crescentemente constrangeram a capacidade de ação
do Estado em termos da sua política fiscal, tributária, de investimentos
e regulação do mercado de trabalho, afetando qualitativa e
quantitativamente a transformação econômica brasileira.

Vale notar que a capacidade de ação do Estado significa mais do que
apenas sua capacidade financeira, apesar desta importar. A ação efetiva
do Estado na economia significa que este consegue atingir objetivos de
acumulação ou distribuição e está ancorada tanto na formação de uma
burocracia eficiente como na obtenção de certa autonomia perante
interesses de classe
\href{https://www.zotero.org/google-docs/?CDUJIb}{(Rueschemeyer e Evans,
1985)}. Também importante para seu sucesso é a relação entre esse
Estado, e suas classes dominantes, e a economia internacional
\href{https://www.zotero.org/google-docs/?HtvVrw}{(Evans, 1985)}.

Assim, longe de meras escolhas conscientes, propõe-se que os arranjos
institucionais a serem analisados foram condicionados pela necessidade
do Estado brasileiro em conciliar interesses: 1) de elites econômicas
nacionais, com lógicas de alocação de capital baseadas no uso de mão de
obra de baixa produtividade ou em atividades rentistas; 2) do capital
estrangeiro, com estratégias de alocação de capital baseados no acesso
ao mercado interno ou de caráter ``extrativista''\footnote{Por
  extrativista, aqui refere-se tanto a atividades tradicionalmente
  denominadas extrativistas, tais como a mineração ou extração de
  petróleo, como também ao uso do poder -- corporativo e institucional
  -- para a extração de valor de atividades econômicas, em oposição a
  estratégias baseadas no incremento do valor produzido
  \href{https://www.zotero.org/google-docs/?GjG8Lj}{(Henderson \emph{et
  al.}, 2002)}. Assim, a alocação de capitais em função de benefícios
  fiscais, baixo custo do trabalho ou leis ambientais menos rigorosas
  poderia ser enquadrada como ``extrativista''.}; 3) das elites
políticas e do funcionalismo público -- notadamente as elites militar,
do judiciário e do ministério público --, que visam projetos ancorados
em interesses próprios, mesmo que por vezes de caráter nacionalista.

Tal cenário se relaciona a um padrão de ``desenvolvimento dependente'',
no qual se observam constrangimentos estruturais para o estado de operar
como agente de transformação econômica e constante instabilidade
derivada de uma inserção precária na economia global
\href{https://www.zotero.org/google-docs/?szdDtj}{(Evans, 1979, 1995)}.
Essa constatação parte do diagnóstico da existência de uma
característica de dependência do desenvolvimento econômico de certos
países no capitalismo, ditos periféricos, em relação aos países
entendidos como centrais
\href{https://www.zotero.org/google-docs/?8DVu6P}{(Cardoso e Faletto,
1970)}. Porém, avança a discussão ao teorizar uma mudança no caráter
dessa dependência em certo conjunto de países a partir de meados do
século XX graças a um intenso processo de industrialização mediado pelo
Estado \href{https://www.zotero.org/google-docs/?slaENo}{(Evans, 1979)}.

Nesse contexto, propõe-se que os arranjos institucionais do caso
brasileiro visam balizar um modelo "conservador" de crescimento dentro
da ótica dessas elites e que responde às instabilidades internas (os
conflitos de interesse das elites) e externas (a relação de
dependência). Exemplos recentes dessas instituições seriam encontradas
na política fiscal de produção de superávit e na independência do Banco
Central do Brasil, por exemplo.

Por fim, a proposta de investigação admite que a desindustrialização,
ocorrida em meio à ``incompletude'' do projeto de desenvolvimento
nacional, mais do que gerar "ruínas"
\href{https://www.zotero.org/google-docs/?mViMNW}{(Mah, 2012)} ou um
estímulo à transição para atividade de serviços intensivos em capital e
de maior rentabilidade
\href{https://www.zotero.org/google-docs/?ISAuAT}{(Amsden, 2003)},
libera mão de obra de modo funcional aos modelos de alocação de capital
de parte das elites nacionais. Assim, a desindustrialização brasileira
representa um rumo distinto dos casos do chamado Norte Global, criando
um caminho de reestruturação da economia em torno de serviços intensivos
em mão de obra por parte das elites nacionais e de estratégias de
alocação de capital de caráter ``extrativista'' por parte do capital
estrangeiro. Nesse cenário, pode coexistir um certo grau de crescimento
econômico com a manutenção de padrões de dependência e trabalho
precarizado.

\subsection{\texorpdfstring{\textbf{2. Justificativa e breve histórico
da desindustrialização
brasileira}}{2. Justificativa e breve histórico da desindustrialização brasileira}}\label{justificativa-e-breve-histuxf3rico-da-desindustrializauxe7uxe3o-brasileira}

Ao longo das últimas décadas, a desindustrialização tem ocorrido em
paralelo às tendências mais gerais de reestruturação produtiva, baseadas
no aumento da produtividade do trabalho industrial e no crescimento
acelerado do setor de serviços
\href{https://www.zotero.org/google-docs/?KKpm5d}{(Rowthorn e Ramaswamy,
1999)}. Entretanto, a variabilidade de seus pontos de partida nos
diversos países
\href{https://www.zotero.org/google-docs/?8KG5nv}{(Rodrik, 2016)} e os
distintos resultados obtidos por estes
\href{https://www.zotero.org/google-docs/?3r6RHy}{(Berger, Musso e
Wicke, 2022)} indicam que instituições de nível nacional e subnacional
são importantes para compreender o grau, ritmo e consequências desse
processo.

A discussão sobre desindustrialização, enquanto um fenômeno global, será
feita de forma mais aprofundada no capítulo 3. Por hora, basta dizer
que, desde meados do século XX, os processos de desindustrialização têm
transformado sociedades em todo mundo
\href{https://www.zotero.org/google-docs/?XIuSoy}{(High, 2013; Lawson,
2020; Rodrik, 2016; Strangleman, 2017; Strangleman, Rhodes e Linkon,
2013)}. Desemprego, austeridade fiscal e um sentimento de desamparo são
algumas das consequências percebidas, com impactos que se perpetuam no
tempo sobre as populações dos territórios afetados
\href{https://www.zotero.org/google-docs/?iOwmqp}{(Linkon, 2018; Mah,
2012)}. Nesses termos, a desindustrialização constitui um dos fenômenos
mais marcantes para os trabalhadores e suas comunidades desde a segunda
metade do século XX.

Processos de desindustrialização como os do Norte da Inglaterra
\href{https://www.zotero.org/google-docs/?U21Pvt}{(Tomlinson, 2016)}, do
\emph{Rust-Belt} americano
\href{https://www.zotero.org/google-docs/?9fBYil}{(High, 2003)} ou mesmo
da Hungria pós-socialismo
\href{https://www.zotero.org/google-docs/?QTEze5}{(Scheiring, 2020)}
exemplificam esses enormes impactos. Outros trabalhos apontam ainda o
aumento da desigualdade
\href{https://www.zotero.org/google-docs/?TXQfOB}{(Kollmeyer, 2018)}, a
piora nas condições de saúde dos trabalhadores
\href{https://www.zotero.org/google-docs/?W7X432}{(Winant, 2019)} e,
contraintuitivamente, a continuidade da operação de atividades nocivas e
da exposição do ambiente e das comunidades a elas
\href{https://www.zotero.org/google-docs/?1xHa2i}{(Feltrin, Mah e Brown,
2022)}\footnote{A atividade ``nociva'' é compreendida aqui como aquela
  que envolve impactos negativos de atividades econômicas para o meio
  ambiente e saúde de trabalhadores. Em particular, ela diz respeito aos
  despejos de resíduos industriais
  \href{https://www.zotero.org/google-docs/?0s96N7}{(Feltrin, Mah e
  Brown, 2022)}.}.

No caso brasileiro, como já dito, a literatura tende a apontar que a
redução do papel da indústria começa a se manifestar de forma mais
intensa ao longo dos anos 1980
\href{https://www.zotero.org/google-docs/?nIqXFH}{(Oreiro e Feijó,
2010)}. Em grande parte, tal retração foi concentrada na indústria de
transformação, de modo que sua participação no PIB nacional se reduziu
de cerca de 27\% para algo próximo de 14\% em fins dos anos 1990
\href{https://www.zotero.org/google-docs/?FecbJJ}{(Morceiro, 2021)}.
Após breve retomada até 18\%, o indicador voltou a se retrair na segunda
metade dos anos 2000, chegando a 11\% em fins da década de 2010 (idem).

Esse processo de desindustrialização de caráter ``prematuro''
\href{https://www.zotero.org/google-docs/?CaOWo7}{(Rodrik,
2016)}\footnote{É considerada ``prematura'' a desindustrialização que
  ocorre em países de renda per capita menos elevada quando comparado ao
  nível de renda per capita dos países onde o fenômeno originalmente se
  manifestou. Nos segundos, a desindustrialização é considerada
  consequência esperada de uma tendência da demanda por serviços crescer
  mais do que a demanda por produtos a partir de dados níveis de renda
  per capita, em paralelo a um crescimento elevado da produtividade
  industrial. O fato do fenômeno ocorrer em níveis de renda per capita
  abaixo desse patamar ``esperado'' configuraria uma desindustrialização
  ``prematura''
  \href{https://www.zotero.org/google-docs/?js82Qa}{(Rodrik, 2016)}.}, é
associado na literatura nacional, majoritariamente de viés econômico, à
perda de competitividade da indústria brasileira em função da abertura
comercial \href{https://www.zotero.org/google-docs/?IgPcre}{(Cano,
2012)}, de sua baixa produtividade
\href{https://www.zotero.org/google-docs/?ETIj2z}{(Oreiro e Marconi,
2012)} e de constantes períodos de sobrevalorização cambial
\href{https://www.zotero.org/google-docs/?1ibNIt}{(Marconi e Rocha,
2012)}, que se relacionam com o redirecionamento do capital nacional
para setores primários da economia
\href{https://www.zotero.org/google-docs/?PD1Cx6}{(Bresser-Pereira,
2020; Oreiro e Feijó, 2010)}. Esses dois últimos elementos, em conjunto,
são entendidos como ``doença holandesa'' ou ``a desindustrialização
causada pela apreciação da taxa real de câmbio resultante da descoberta
de recursos naturais escassos num determinado país ou região''
\href{https://www.zotero.org/google-docs/?7SQuPc}{(Oreiro e Feijó, 2010,
p. 222)}. São, entretanto, pouco aprofundadas nessa literatura as
discussões sobre o impacto das instituições políticas e do caráter
``dependente'' do desenvolvimento brasileiro
\href{https://www.zotero.org/google-docs/?1XbSAc}{(Evans, 1979)}, que
condicionam as possibilidades e limites do processo de acumulação de
capital na periferia do capitalismo
\href{https://www.zotero.org/google-docs/?Bfmt8N}{(Evans, 1995)}.

Para alguns autores, essa redução da participação da indústria na
economia é particularmente relevante pois afeta o potencial de
crescimento econômico brasileiro
\href{https://www.zotero.org/google-docs/?dCuzii}{(Oreiro e Marconi,
2012)}. Além disso, o setor industrial nacional tipicamente oferecia
empregos com maior estabilidade, maiores salários, mais proteção social
e maiores taxas de sindicalização que outros setores
\href{https://www.zotero.org/google-docs/?8bSstM}{(Ramalho, 2004)}.
Mesmo as novas unidades produtivas inauguradas desde os anos 1990, que
reproduzem de maneira incompleta a qualidade na geração de postos de
trabalho, por já nascerem em um contexto de ``flexibilização", ainda
oferecem empregos significativamente melhores -- nos termos
supramencionados -- e capazes de impulsionar melhorias substanciais em
territórios \href{https://www.zotero.org/google-docs/?UbZXip}{(Ramalho,
2015)}.

Assim, a desindustrialização representa, no caso brasileiro, uma redução
dos horizontes de possibilidade de melhoria econômica e social. Isso se
agrava pelo fato de que o setor de serviços brasileiro compreende, à
exceção de poucos casos, atividades economicamente desvalorizadas ou
caracterizadas por piores condições de trabalho. Dessa forma, dinâmicas
de reestruturação da economia em torno do setor de serviços, que estão
embutidas em processos de desindustrialização, não remetem, ao menos
prioritariamente, à migração de profissionais para setores de tecnologia
ou finanças e de ampliação de postos de trabalho que demandam
conhecimentos. No Brasil e mesmo na América Latina mais amplamente, se
trata da expansão de serviços precários, com ênfase em serviços
``intensivos'' em mão de obra como \emph{call centers} ou, mais
recentemente , trabalho em plataformas digitais
\href{https://www.zotero.org/google-docs/?lOg0a8}{(Bridi, Oliveira e
Salas, 2023)}.

Dada a desigualdade e a diversidade dos territórios brasileiros,
entretanto, é possível que os efeitos desse da desindustrialização não
sejam igualmente visíveis em todo país. Mesmo num contexto de
desindustrialização em escala nacional, um conjunto de estados
brasileiros têm expandido a participação do setor secundário em sua
matriz econômica, assim como vem diversificando sua estrutura
industrial, em meio ao que pode ser entendido como uma dinâmica de
desconcentração industrial do Brasil
\href{https://www.zotero.org/google-docs/?mO9EhC}{(CNI, 2021)}. Como os
efeitos de intensa desindustrialização em dado território podem ser
parcialmente mascarados pelo crescimento industrial em outros, é
relevante se concentrar nos estados e territórios que melhor representem
o fenômeno. Os dois casos que mais se adequam a essa condição são,
assim, os estados de São Paulo e do Rio de Janeiro.

O estado de São Paulo possui o maior PIB industrial do Brasil
\href{https://www.zotero.org/google-docs/?Up7fFY}{(CNI, {[}s.d.{]})},
contribuindo com 26\% da indústria nacional. Em relação à economia
estadual, a indústria representa 23\% desta e emprega 3.3 milhões de
pessoas, cerca de 22\% do emprego formal do estado, com base nos dados
mais atuais \href{https://www.zotero.org/google-docs/?NASjAt}{(CNI,
{[}s.d.{]})}. Mesmo diante do processo de desconcentração industrial
supracitado, o estado de São Paulo ainda constitui o elo mais importante
da dinâmica industrial brasileira.

Essa liderança está associada à longa importância do estado para a
economia brasileira, que data da economia cafeeira de fins do século
XIX. Esta, ancorada na escravidão, provocou dinâmicas de acumulação e
realocação de capital capazes de impulsionar o processo de
industrialização do estado
\href{https://www.zotero.org/google-docs/?dqEUDA}{(Lanna, 2019)},
capitais em grande parte disponíveis após a abolição da escravidão .

Em função da sua centralidade econômica, o estado também ficou marcado
por significativos fluxos de imigrantes, em parte estrangeiros e em
parte de outras regiões brasileiras
\href{https://www.zotero.org/google-docs/?Q9aFbc}{(Lanza e Lamounier,
2014)}. Central nessa dinâmica foi a relação estabelecida entre São
Paulo, como \emph{locus} da industrialização nacional, e o Nordeste
brasileiro, enquanto origem de grande parte de sua força de trabalho,
especialmente após a segunda guerra mundial
\href{https://www.zotero.org/google-docs/?MI4Xs4}{(Fontes, 2008)}. Essa
``tradição industrial'' se perpetuou durante todo século XX e foi
perpassada por diversos projetos nacionais de ``modernização'' ou
desenvolvimento. Em particular, se ligará com o modelo rodoviarista, na
medida em que se constituiu também no principal mercado consumidor
brasileiro, obtendo do setor automotivo significativo impulso adicional
à sua industrialização
\href{https://www.zotero.org/google-docs/?I30f5z}{(Dulci, 2021)}.

Talvez em nenhuma outra parte a expressão desse modelo rodoviarista e do
impulso da indústria automotiva seja tão visível quanto na sub-região do
ABC Paulista (doravante apenas ABC). O ABC representa um conjunto de
municípios localizados ao sul da cidade de São Paulo, ainda em sua
região metropolitana. A cidade de São Paulo é localizada ao Norte do
porto de Santos, importante centro de entrada de escravizados e
escoamento de café no século XIX
\href{https://www.zotero.org/google-docs/?yNmXPa}{(Lanna, 2019)}, e
integra hoje a principal aglomeração urbana brasileira, a Região
Metropolitana de São Paulo (RMSP), que conta com cerca de 21 milhões de
habitantes \href{https://www.zotero.org/google-docs/?SCF7Gb}{(Governo do
Estado de São Paulo, {[}s.d.{]})} Na RMSP, os municípios que compõem o
ABC somam 2,7 milhões de habitantes
\href{https://www.zotero.org/google-docs/?hUGjhR}{(Consórcio ABC, 2024)}
e constituem um denso território produtivo, em grande parte associado a
empreendimentos da indústria automotiva, como já sugerido.

Mas a importância histórica do ABC vai além do aspecto econômico. Em
grande parte, ela é atrelada aos trabalhadores do setor automotivo, de
modo que a sub-região viu o surgimento de um significativo movimento
sindical relacionado a essa atividade. Esse movimento exerceu liderança
nas disputas sobre a pauta trabalhista, em especial nos anos 1970 e 1980
\href{https://www.zotero.org/google-docs/?6z3SJT}{(Ramalho e Rodrigues,
2013, 2018)}, de modo que movimentos grevistas da sub-região foram, por
exemplo, importante elemento na luta pela redemocratização brasileira.
Essa liderança teve como consequência a formação de líderes políticos,
como o presidente Luiz Inácio Lula da Silva, que impactaram
significativamente a história recente do Brasil
\href{https://www.zotero.org/google-docs/?do1NWN}{(Ramalho e Rodrigues,
2018)}. Nesse sentido, a desindustrialização de São Paulo e do ABC em
particular, representa também uma reestruturação da classe trabalhadora
e um significativo impacto sobre sua capacidade de mobilização e
reivindicação \href{https://www.zotero.org/google-docs/?X2RzFf}{(Fontes,
2023)}.

Daí também a importância da reestruturação produtiva, que se intensifica
nos anos 1990 \href{https://www.zotero.org/google-docs/?DEVtfA}{(Ramalho
e Rodrigues, 2013)}. Como resultado, vemos uma constante redução da mão
de obra industrial, em termos relativos, em São Paulo desde os anos 1980
\href{https://www.zotero.org/google-docs/?O9uEV4}{(Sobral, 2016)}. Ao
longo dos anos 2000 e 2010 as tendências de desindustrialização se
mantêm. Apenas entre 2011 e 2021 o estado perdeu 4,2 pontos percentuais
em sua participação no PIB industrial nacional
\href{https://www.zotero.org/google-docs/?2YaMK0}{(CNI, {[}s.d.{]})}. Ao
observar os dados de emprego no setor de transformação por estado,
apesar de melhor do que a situação do Rio de Janeiro, fica claro que a
participação do emprego no setor se reduziu em relação a outros estados
do país \href{https://www.zotero.org/google-docs/?Xjf771}{(Sobral,
2016)} (cf. tabela 1 abaixo). No caso da sub-região do ABC, os dados
mostram ainda uma redução do número de empresas desde 1998
\href{https://www.zotero.org/google-docs/?D25gWL}{(SMABC, 2004)} e uma
aceleração da perda do seu PIB industrial devido a disputas fiscais
entre estados, muito acentuada durante o período de 2013 até 2016
\href{https://www.zotero.org/google-docs/?BrK8oO}{(Valor Econômico,
2019)}.

Em função desses impactos políticos e econômicos, tentativas de reação
foram articuladas
\href{https://www.zotero.org/google-docs/?rK1JZk}{(Reis, 2007)}.
Destaque pode ser conferido aos esforços de construção institucional
intermunicipais, no que ficou consolidado como Consórcio do ABC
\href{https://www.zotero.org/google-docs/?88O0x4}{(Ramalho e Rodrigues,
2010, 2013)}. O fato de um conjunto de lideranças políticas regionais
estar diretamente ligado ao Partido dos Trabalhadores, que durante a
maior parte dos anos 2010 e 2010 também ocupava a presidência da
república, permitiu a articulação de projetos de reestruturação
econômica em torno da tentativa de manutenção dos empreendimentos
industriais. Exemplo disso é o caso da tentativa de reconversão
industrial orientada para o setor de defesa
\href{https://www.zotero.org/google-docs/?2jRDZr}{(Ramalho e Conceição,
2024)}.

Tabela 1: Empregos no setor de transformação entre 1985 e 2012 em
estados selecionados

\begin{longtable}[]{@{}
  >{\raggedright\arraybackslash}p{(\linewidth - 6\tabcolsep) * \real{0.4991}}
  >{\raggedright\arraybackslash}p{(\linewidth - 6\tabcolsep) * \real{0.1809}}
  >{\raggedright\arraybackslash}p{(\linewidth - 6\tabcolsep) * \real{0.1635}}
  >{\raggedright\arraybackslash}p{(\linewidth - 6\tabcolsep) * \real{0.1565}}@{}}
\toprule\noalign{}
\begin{minipage}[b]{\linewidth}\raggedright
\textbf{Estados}
\end{minipage} & \begin{minipage}[b]{\linewidth}\raggedright
\textbf{1985}
\end{minipage} & \begin{minipage}[b]{\linewidth}\raggedright
\textbf{2012}
\end{minipage} & \begin{minipage}[b]{\linewidth}\raggedright
\textbf{Variação}
\end{minipage} \\
\begin{minipage}[b]{\linewidth}\raggedright
São Paulo
\end{minipage} & \begin{minipage}[b]{\linewidth}\raggedright
2492802
\end{minipage} & \begin{minipage}[b]{\linewidth}\raggedright
2820813
\end{minipage} & \begin{minipage}[b]{\linewidth}\raggedright
13.16\%
\end{minipage} \\
\begin{minipage}[b]{\linewidth}\raggedright
Rio de Janeiro
\end{minipage} & \begin{minipage}[b]{\linewidth}\raggedright
520334
\end{minipage} & \begin{minipage}[b]{\linewidth}\raggedright
464796
\end{minipage} & \begin{minipage}[b]{\linewidth}\raggedright
-10.67\%
\end{minipage} \\
\begin{minipage}[b]{\linewidth}\raggedright
Rio Grande do Sul
\end{minipage} & \begin{minipage}[b]{\linewidth}\raggedright
502318
\end{minipage} & \begin{minipage}[b]{\linewidth}\raggedright
733387
\end{minipage} & \begin{minipage}[b]{\linewidth}\raggedright
46.00\%
\end{minipage} \\
\begin{minipage}[b]{\linewidth}\raggedright
Minas Gerais
\end{minipage} & \begin{minipage}[b]{\linewidth}\raggedright
392529
\end{minipage} & \begin{minipage}[b]{\linewidth}\raggedright
841694
\end{minipage} & \begin{minipage}[b]{\linewidth}\raggedright
114.43\%
\end{minipage} \\
\begin{minipage}[b]{\linewidth}\raggedright
Santa Catarina
\end{minipage} & \begin{minipage}[b]{\linewidth}\raggedright
286345
\end{minipage} & \begin{minipage}[b]{\linewidth}\raggedright
641212
\end{minipage} & \begin{minipage}[b]{\linewidth}\raggedright
123.93\%
\end{minipage} \\
\begin{minipage}[b]{\linewidth}\raggedright
Paraná
\end{minipage} & \begin{minipage}[b]{\linewidth}\raggedright
235514
\end{minipage} & \begin{minipage}[b]{\linewidth}\raggedright
678080
\end{minipage} & \begin{minipage}[b]{\linewidth}\raggedright
187.91\%
\end{minipage} \\
\begin{minipage}[b]{\linewidth}\raggedright
\textbf{Média dos 6 estados}
\end{minipage} & \begin{minipage}[b]{\linewidth}\raggedright
\textbf{738307}
\end{minipage} & \begin{minipage}[b]{\linewidth}\raggedright
\textbf{1029997}
\end{minipage} & \begin{minipage}[b]{\linewidth}\raggedright
\textbf{39.51\%}
\end{minipage} \\
\multicolumn{4}{@{}>{\raggedright\arraybackslash}p{(\linewidth - 6\tabcolsep) * \real{1.0000} + 6\tabcolsep}@{}}{%
\begin{minipage}[b]{\linewidth}\raggedright
Fonte: adaptado de Sobral
\href{https://www.zotero.org/google-docs/?NGRu07}{(2016, p. 198)}, com
base na RAIS/MTE
\end{minipage}} \\
\midrule\noalign{}
\endhead
\bottomrule\noalign{}
\endlastfoot
\end{longtable}

O que essa mobilização demonstra, prioritariamente, é o forte caráter
político do processo de reestruturação, que envolve projetos e arranjos
sociopolíticos voltados a influenciar a dinâmica econômica regional.
Adicionalmente, essa mobilização encontra recurso simbólico importante
na construção de uma identidade comum aos municípios do ABC, que acaba
justamente sendo reforçada durante momentos de crise. É assim, um caso
privilegiado para uma compreensão multidimensional -- política,
econômica e cultural -- do processo de desindustrialização e também das
suas relações com transformações institucionais.

Já o estado do Rio de Janeiro representa uma antiga tradição industrial
que, ao longo de todo século XX, se deparou com ciclos de crise e
gradual declínio. De fato, a atividade industrial no estado é talvez a
mais antiga no país, remontando à chegada da Família Real e ao
estabelecimento da capital do império no Rio de Janeiro. Pouco antes da
Independência (1822) e estabelecida como província em 1821, ainda por
muito tempo o Rio de Janeiro foi uma peça central da atividade econômica
nacional devido a sua produção cafeeira. Nesse contexto, os esforços de
industrialização eram majoritariamente ligados aos esforços da coroa em
constituir capacidade produtiva para atender uma crescente demanda,
excedentes produtivos advindos da agricultura e eventuais tentativas de
substituição de importações
\href{https://www.zotero.org/google-docs/?r9UnQR}{(Barreto, 2018)}

De fins do século XIX até o início do século XX, entretanto, o eixo da
atividade cafeeira se deslocou do Rio de Janeiro para São Paulo, com a
crise do Vale do Paraíba. Devido, em grande parte, à manutenção do
distrito federal no Rio de Janeiro, desde 1891 um estado, este se
manteve até princípios do século XX como maior polo industrial. Mas, até
1930, São Paulo assumiria a liderança industrial
\href{https://www.zotero.org/google-docs/?mkI5eI}{(Saes e Nozoe, 2014)}.
De toda maneira, a importância do estado do Rio de Janeiro permanecia e
este se mantinha como a segunda potência industrial do país e também
casa da capital federal.

É nesse contexto que, em 1941, seria estabelecida a Companhia
Siderúrgica Nacional (CSN), importante evento para a cronologia deste
trabalho. Criada por decreto de Vargas em meio a negociações sobre a
entrada do Brasil na segunda guerra mundial, a CSN constituirá uma peça
fundamental do processo de industrialização por substituição de
importações \href{https://www.zotero.org/google-docs/?mSOr4o}{(Ramalho,
Santos e Lima, 2013)}. Sua instalação é em Volta Redonda, o principal
centro urbano da sub-região do Sul fluminense.

Em torno da instalação dessa siderúrgica houve não apenas a
significativa expansão urbana, mas também a formação de uma classe
trabalhadora industrial
\href{https://www.zotero.org/google-docs/?3VJEnI}{(Morel, 1989)}. Entre
os anos 1940 e 1970, a promoção estatal da instalação de empreendimentos
industriais, a exemplo da Fábrica Nacional de Motores (FNM)
\href{https://www.zotero.org/google-docs/?YHTmgi}{(Ramalho, 1989)}, e a
própria dinâmica econômica do estado preservaram um tecido industrial
ainda expressivo. Assim, o território produtivo do Sul fluminense se
consolidou como tradicionalmente industrial.

Entretanto, já nos anos 1980 e diante da crise da dívida enfrentada pelo
país, a indústria fluminense passaria a observar uma retração bastante
rápida das unidades industriais\footnote{Santos
  \href{https://www.zotero.org/google-docs/?1lM2aG}{(2006)} traça um
  quadro do encerramento de atividades industriais importantes no Sul do
  estado do Rio de Janeiro nos anos 1980 e 1990. Mas essa dinâmica
  possui feição mais ampla, afetando todo o Rio de Janeiro.}. No caso da
CSN e de Volta Redonda, em particular, a gradual retração do número de
trabalhadores siderúrgicos e a perda de importância da siderúrgica em
âmbito nacional contribuiu de maneira significativa para a crise da
sub-região \href{https://www.zotero.org/google-docs/?5M72z0}{(Ramalho,
Santos e Lima, 2013)}.

A resposta a tal crise foi observada na forma de arranjos sociopolíticos
orientados a atrair empresas do setor automotivo
\href{https://www.zotero.org/google-docs/?zKdBOM}{(Ramalho, 2015;
Santos, 2007)}. O principal marco do processo diz respeito à instalação
do Consórcio Modular da Volkswagen em Resende
\href{https://www.zotero.org/google-docs/?tDsDc8}{(Abreu, Beynon e
Ramalho, 2000)}. Posteriormente, uma série de investimentos consolidou a
sub-região como parte de uma rede global de produção automotiva
\href{https://www.zotero.org/google-docs/?lUPNv9}{(Santos, 2021)}.

Em que pesem as limitadas evidências de que tais investimentos
impulsionaram a atividade industrial de outros setores, fica claro que o
Sul fluminense experimentou a reversão de sua perda de capacidade
produtiva industrial, de postos de trabalho e, na esteira destas,
verificou melhorias nas condições sociais, amparada na capacidade fiscal
ampliada dos municípios
\href{https://www.zotero.org/google-docs/?YKljwf}{(Ramalho, 2015)}. O
processo também serviu de base, especialmente em períodos de crise, para
a articulação de novos arranjos institucionais que visavam remediar
crises econômicas e setoriais, como o Fórum Demissão Zero
\href{https://www.zotero.org/google-docs/?kzIW7G}{(Ramalho, 2005)}.
Assim, a sub-região representa, mesmo que com pouca participação no PIB
fluminense e na indústria nacional, um caso privilegiado para o estudo
das relações entre transformações institucionais e processos de
desindustrialização ou reindustrialização.

Apesar disso, a recuperação da dinâmica industrial da sub-região a
partir dos anos 90 não se refletia no estado como um todo. Isso mesmo
com a implantação da indústria petroquímica no Norte do estado
contribuindo de modo decisivo para a desaceleração da retração
industrial \href{https://www.zotero.org/google-docs/?xIxk42}{(Piquet,
2021)}. Além de ver ameaçada por Minas Gerais sua posição como segunda
maior potência industrial
\href{https://www.zotero.org/google-docs/?uC5UHk}{(CNI, 2021)}, o Rio de
Janeiro gradativamente perdeu espaço na produção nacional para estados
como Bahia, Pernambuco, Paraná, Rio Grande do Sul e Mato Grosso do Sul
\href{https://www.zotero.org/google-docs/?TtN18w}{(CNI, 2021)}.

Outro importante marco da perda de centralidade econômica e política do
Rio de Janeiro ao longo do século XX é a transferência da capital para
Brasília (1960). Agora distante do núcleo decisório da política nacional
e progressivamente menos integrado aos projetos de desenvolvimento ali
estabelecidos, o Rio de Janeiro vê sua importância econômica decrescer
frente ao resto do país desde então.

É importante também destacar os efeitos de todo esse processo para as
cidades que compõem a Região Metropolitana do Rio de Janeiro (RMRJ),
onde a tendência à concentração da atividade econômica na capital
ensejou a formação de vazios produtivos, com forte consequência para os
municípios da Baixada Fluminense
\href{https://www.zotero.org/google-docs/?cXsZIX}{(Sobral, 2016, 2017)}.
Ao mesmo tempo, a própria capital sofria com o encerramento de fábricas
e a redução de empregos industriais, em meio a uma transformação de seu
mercado de trabalho, cada vez mais orientado ao setor de serviços .

Um conjunto de dados pode ajudar a situar tal cenário. Como mostra
Sobral \href{https://www.zotero.org/google-docs/?gq3DiJ}{(2017)}, o
estado do Rio de Janeiro teve um crescimento de sua produção física
industrial significativamente abaixo de São Paulo entre 1995 e 2010. O
mesmo autor aponta uma estagnação na produção física da indústria de
transformação fluminense no mesmo período, em contraste com um
crescimento de cerca de 40\% tanto na média brasileira como em São Paulo
no mesmo período e cerca de 45\% e 50\% nos casos de estados como Bahia
e Rio Grande do Sul
\href{https://www.zotero.org/google-docs/?cytWuo}{(Sobral, 2017, p.
11)}.

Dados sobre o período mais recente, de 2011 até 2021, mostram uma
pequena oscilação positiva da participação da indústria no PIB estadual
e da participação do estado no PIB industrial nacional
\href{https://www.zotero.org/google-docs/?PQBOvv}{(CNI, {[}s.d.{]})},
mas os dados das Contas Regionais
\href{https://www.zotero.org/google-docs/?qrmnq0}{(IBGE, {[}s.d.{]})}
mostram que tal variação está ligada quase exclusivamente ao setor
extrativo. Dados sobre a performance dos subsetores quando comparados ao
PIB do estado reforçam que a indústria de transformação tem um
desempenho consistentemente inferior ao do PIB estadual, enquanto os
setores extrativo e de construção oscilam (cf. figura 1 abaixo). Tal
dado é agravado pela redução constante da mão de obra industrial
\href{https://www.zotero.org/google-docs/?I3Hbh4}{(Sobral, 2016, p.
198)} (tabela 1, p.14) e pela manutenção de taxas de crescimento
industrial abaixo do nível brasileiro. De fato, toda a performance da
economia fluminense parece estar abaixo dos níveis nacionais,
sinalizando o longo desenrolar de uma crise que é tanto política como
econômica.

Figura 1: taxa de crescimento real dos subsetores industriais comparados
ao PIB estadual

\includegraphics[width=6.5in,height=3.31944in]{media/image1.png}

Fonte: Autoria própria a partir das Contas Regionais
\href{https://www.zotero.org/google-docs/?PfGVry}{(IBGE, {[}s.d.{]})}

\subsection{\texorpdfstring{\textbf{3 - Instituições, comparações
históricas e desenvolvimento
dependente}}{3 - Instituições, comparações históricas e desenvolvimento dependente}}\label{instituiuxe7uxf5es-comparauxe7uxf5es-histuxf3ricas-e-desenvolvimento-dependente}

Diferente da teoria substantiva, que versa sobre as especificidades do
problema investigado, a teoria geral, aqui apresentada ainda de forma
preliminar, fornece o enquadramento mais abstrato dentro do qual se
analisa o fenômeno
\href{https://www.zotero.org/google-docs/?YuggLP}{(Sautu \emph{et al.},
2005)}. Nesse trabalho, ela abarca a vertente institucionalista das
análises históricas e comparativas nas ciências sociais, que por vezes
podem ser encontradas como parte de campos maiores de diferentes
denominações como: a análise histórico comparativa
\href{https://www.zotero.org/google-docs/?apTiPh}{(Mahoney e
Rueschemeyer, 2003; Thelen e Mahoney, 2015)}, a sociologia histórica
\href{https://www.zotero.org/google-docs/?OrkRyO}{(Skocpol, 1984a)}, a
análise institucionalista comparada
\href{https://www.zotero.org/google-docs/?fkkDeO}{(Thelen, 2010)} ou a
sociologia econômica
\href{https://www.zotero.org/google-docs/?cd8zSi}{(Dobbin, 2005)}.

Tal vertente é centrada na premissa de que processos institucionalização
e transformação institucional são parte central da vida social e
precisam ser avaliados enquanto processos ocorrendo no tempo e não
recortes estáticos. Essas perspectivas são postas em diálogo com as já
abordadas teorias sobre o papel do Estado nos processos de transformação
econômica \href{https://www.zotero.org/google-docs/?Yfh3Nf}{(Evans,
Rueschemeyer e Skocpol, 1985)}, particularmente a discussão sobre
desenvolvimento dependente
\href{https://www.zotero.org/google-docs/?DYrm0Y}{(Evans, 1979)}.

Em que pese múltiplas e frequentes disputas sobre o conceito de
instituição \href{https://www.zotero.org/google-docs/?QiRe92}{(Duina,
2011; Morgan \emph{et al.}, 2010; Nee, 2005; Scott, 2014)}, autores
convergem na sua caracterização como padrão orientador da ação, sendo
percebida como conjunto de regras com níveis variados de formalização e
podendo se constituir de elementos materiais e imateriais. Com o
objetivo de ampliar sua clareza, a tese adota uma definição de
instituições como sistemas de crenças, normas e expectativas socialmente
compartilhadas que assumem caráter de regras mais ou menos formalizadas
e que enquadram, habilitam ou constrangem a forma como atores de
diversos níveis percebem o mundo e atuam nele.

Adicionalmente, dentro da sociologia econômica, especialmente a de
caráter histórico e comparativo
\href{https://www.zotero.org/google-docs/?Mnc7j7}{(Dobbin, 2005)},
perspectivas institucionais são complementadas com discussão sobre redes
e poder. Quanto às primeiras, autores nessa subárea mobilizam com
frequência a ideia de redes, enfocando as conexões -- formais ou
informais -- que atores estabelecem entre si e através das quais obtêm
informações e/ou recursos ou exercem influência e pressão. A inserção de
atores em redes muitas vezes condiciona sua ação, o que Granovetter
\href{https://www.zotero.org/google-docs/?CG4uYo}{(2007)} irá conceituar
como ``enraizamento''. Esse conceito, porém, pode ser expandido para
além de seu papel nas redes, incluindo também dimensões societal e
territorial \href{https://www.zotero.org/google-docs/?syvXQo}{(Hess,
2004)}. Já as relações de poder, em grande parte imersas nas redes e
instituições, são a capacidade de certos atores promoverem seus próprios
interesses, em especial sua habilidade de promovê-los enquanto supostos
interesses comuns
\href{https://www.zotero.org/google-docs/?j9POT7}{(Dobbin, 2005)}.

Por sua vez, nas análises de cunho histórico comparativo, como a aqui
proposta, algumas importantes ferramentas analíticas são mobilizadas.
Destaca-se aqui os conceitos de trajetória de dependência e de
conjuntura crítica
\href{https://www.zotero.org/google-docs/?cnzUw0}{(Mahoney, 2000,
2004)}, que operacionalizam a abordagem processual pretendida.

A dependência de trajetória
\href{https://www.zotero.org/google-docs/?iwUinW}{(Mahoney, 2000)} se
refere aos efeitos recursivos que reforçam escolhas passadas,
contribuindo para certa estabilidade e prolongando os impactos de
escolhas feitas em momentos-chave. Nesse sentido, a trajetória
dependente é o que constitui a relativa estabilidade dos arranjos
institucionais à medida que estes são constituídos.

Já a conjuntura crítica
\href{https://www.zotero.org/google-docs/?Tg9wjS}{(Mahoney, 2004)} diz
respeito aos momentos chave em que tais escolhas podem ser feitas -- ou
não -- e que têm potencial de alterar essas trajetórias, muitas vezes em
função de um choque ou fator exógeno. Diferentemente da dependência de
trajetória, a conjuntura crítica responde por mudanças substanciais de
rotas institucionais. É importante destacar que isso não implica que
instituições se modifiquem apenas durante conjunturas críticas, como
mostra Thelen \href{https://www.zotero.org/google-docs/?fJR3cT}{(Thelen,
2010)}. Na verdade, as transformações graduais são frequentemente
empregadas por atores que se aproveitam de brechas nas regras ou fazem
uso de poder para direcionar ou recriar instituições a partir de seus
interesses. Estas oportunidades apenas se tornam mais claras em
conjunturas críticas, que tendem a ter mudanças de curto prazo mais
dramáticas.

Assim, pode-se pensar esquematicamente como os três elementos se
relacionam. Redes são meios fundamentais para o exercício do poder. Este
por sua vez é mobilizado por atores na tentativa de construir mundos
estáveis nos quais seus interesses são vistos como comuns e as
incertezas advindas das disputas podem ser reduzidas
\href{https://www.zotero.org/google-docs/?IQxn6o}{(Fligstein, 1996)}. As
instituições cumprem esse papel, ao mesmo tempo que reforçam
recursivamente a dimensão de poder dos atores, conferindo-lhe recursos e
legitimidade.

Outra chave importante para a discussão diz respeito à análise do
processo de desenvolvimento econômico do Brasil no âmbito de um modelo
de desenvolvimento ``dependente''
\href{https://www.zotero.org/google-docs/?uIGiBo}{(Evans, 1979, 2023)}.
Para esse modelo, mesmo que a ação do Estado contribua para um intenso
processo de industrialização em alguns países da periferia do
capitalismo, a combinação da estrutura de classe desses países e com sua
relação com as economias centrais do capitalismo resulta na manutenção
de uma situação de dependência. Essa avaliação tem como base a análise
da dinâmica entre dependência de capital externo, na forma de
multinacionais, uma burguesia nacional com capacidade de investimento e
competição restrita e um Estado que atua como mediador desses interesses
e buscando o estabelecimento e expansão de setores estratégicos. O
resultado são processos de industrialização que convivem com
significativa exclusão social e, por diversas vezes, autoritarismo
\href{https://www.zotero.org/google-docs/?ecY6hE}{(Evans, 1979)}.

Evans \href{https://www.zotero.org/google-docs/?9sf4J2}{(1995)} vai,
entretanto, ampliar a discussão, pensando nas diferentes possibilidades
pelas quais a ação estatal será indutora de processos de
industrialização, mesmo nesse cenário de contingência. Dialoga portanto
com uma retomada da importância dada ao papel do Estado nas ciências
sociais \href{https://www.zotero.org/google-docs/?nUt3EE}{(Evans,
Rueschemeyer e Skocpol, 1985)}\footnote{Esse debate se vincula à
  discussão sobre as condições diferenciais para a industrialização,
  segmentando experiências originária (Inglaterra), tardia (Alemanha e
  EUA, por exemplo) e retardatária (Brasil e Japão, por exemplo),
  conforme periodização proposta por Cardoso de Mello
  \href{https://www.zotero.org/google-docs/?xGNYgb}{(1982)}. Enquanto na
  primeira experiência o papel do Estado é pouco relevante, no
  capitalismo tardio, a escala da atividade econômica demandaria a ação
  estatal no âmbito do financiamento, como indica Gerschenkron
  \href{https://www.zotero.org/google-docs/?THn7ON}{(1962)}. Os desafios
  à industrialização seriam ainda maiores no âmbito da industrialização
  retardatária, exigindo a ação do Estado não apenas quanto ao
  financiamento, mas particularmente quanto à identificação das
  oportunidades de investimento e seu aproveitamento efetivo, como a
  análise de Hirschman
  \href{https://www.zotero.org/google-docs/?iB2ElZ}{(1958)} aponta.},
que também pode ser vista nas perspectivas sobre capitalismo de Estado
\href{https://www.zotero.org/google-docs/?jdQydT}{(Allen, Wood e Keller,
2022)} e no reconhecimento da dimensão política que cerca a construção
de mercados e as estratégias corporativas dentro destes
\href{https://www.zotero.org/google-docs/?WDxVMO}{(Fligstein, 1996)}.

\subsection{\texorpdfstring{\textbf{4.
Objetivos}}{4. Objetivos}}\label{objetivos}

O objetivo geral da tese é compreender como instituições políticas e
econômicas condicionam processos de transformação econômica através do
caso da desindustrialização brasileira. Dentro deste, alguns objetivos
específicos podem ser traçados.

\begin{itemize}
\item
  Discutir a teoria institucionalista de viés histórico comparativo e a
  respeito do papel do Estado na economia, especialmente em relação a
  casos de ``desenvolvimento dependente''
\item
  Contribuir com desenvolvimento de métodos mistos nas ciências sociais,
  particularmente combinando comparação histórica com análise de séries
  temporais e \emph{process tracing}
\item
  Discutir o fenômeno da desindustrialização em escala global, pensando
  suas principais características e sua variabilidade histórica e
  geográfica
\item
  Identificar as transformações institucionais que mais afetaram o
  processo de desindustrialização brasileiro e seus impactos sobre
  diferentes subsetores da indústria e da economia brasileira em geral
\item
  Identificar similaridades e diferenças do processo de
  desindustrialização brasileiro na escala estadual e o impacto de
  instituições de nível subnacional
\item
  Compreender o papel de estratégias corporativas de realocação
  produtiva dentro de processos mais amplos de desindustrialização e
  como arranjos sociopolíticos no nível sub-regional respondem a essas
  pressões
\item
  Integrar a perspectiva multiescalar, produzindo uma síntese que
  contribua para compreensão da especificidade do caso brasileiro
\end{itemize}
